% Options for packages loaded elsewhere
\PassOptionsToPackage{unicode}{hyperref}
\PassOptionsToPackage{hyphens}{url}
\PassOptionsToPackage{dvipsnames,svgnames,x11names}{xcolor}
\documentclass[
  11pt,
]{article}
\usepackage{xcolor}
\usepackage[margin=2.2cm]{geometry}
\usepackage{amsmath,amssymb}
\setcounter{secnumdepth}{-\maxdimen} % remove section numbering
\usepackage{iftex}
\ifPDFTeX
  \usepackage[T1]{fontenc}
  \usepackage[utf8]{inputenc}
  \usepackage{textcomp} % provide euro and other symbols
\else % if luatex or xetex
  \usepackage{unicode-math} % this also loads fontspec
  \defaultfontfeatures{Scale=MatchLowercase}
  \defaultfontfeatures[\rmfamily]{Ligatures=TeX,Scale=1}
\fi
\usepackage{lmodern}
\ifPDFTeX\else
  % xetex/luatex font selection
\fi
% Use upquote if available, for straight quotes in verbatim environments
\IfFileExists{upquote.sty}{\usepackage{upquote}}{}
\IfFileExists{microtype.sty}{% use microtype if available
  \usepackage[]{microtype}
  \UseMicrotypeSet[protrusion]{basicmath} % disable protrusion for tt fonts
}{}
\makeatletter
\@ifundefined{KOMAClassName}{% if non-KOMA class
  \IfFileExists{parskip.sty}{%
    \usepackage{parskip}
  }{% else
    \setlength{\parindent}{0pt}
    \setlength{\parskip}{6pt plus 2pt minus 1pt}}
}{% if KOMA class
  \KOMAoptions{parskip=half}}
\makeatother
\usepackage{longtable,booktabs,array}
\newcounter{none} % for unnumbered tables
\usepackage{calc} % for calculating minipage widths
% Correct order of tables after \paragraph or \subparagraph
\usepackage{etoolbox}
\makeatletter
\patchcmd\longtable{\par}{\if@noskipsec\mbox{}\fi\par}{}{}
\makeatother
% Allow footnotes in longtable head/foot
\IfFileExists{footnotehyper.sty}{\usepackage{footnotehyper}}{\usepackage{footnote}}
\makesavenoteenv{longtable}
\setlength{\emergencystretch}{3em} % prevent overfull lines
\providecommand{\tightlist}{%
  \setlength{\itemsep}{0pt}\setlength{\parskip}{0pt}}

\usepackage{fontspec}
% Symbols missing from the text fonts (arrows, maths signs) come from these fallbacks.
\directlua{luaotfload.add_fallback("symbols", {"STIX Two Math:mode=harf;", "DejaVu Sans:mode=harf;"})}
\setmainfont{STIX Two Text}[RawFeature={fallback=symbols}]
\setmonofont{DejaVu Sans Mono}[Scale=0.85, RawFeature={fallback=symbols}]
\usepackage{unicode-math}
\setmathfont{STIX Two Math}
\usepackage{tikz}
\usetikzlibrary{graphs, graphdrawing, quotes, arrows.meta}
\usegdlibrary{layered}
\usepackage{adjustbox}
\usepackage{float}
\usepackage{fvextra}
\fvset{breaklines=true, breakanywhere=true}
\RecustomVerbatimEnvironment{verbatim}{Verbatim}{breaklines=true, breakanywhere=true}
\renewcommand{\arraystretch}{1.2}
% Let identifiers like learn_temperature=True break after underscores (e.g. in table cells).
\let\oldunderscore\_
\renewcommand{\_}{\oldunderscore\allowbreak}
\usepackage{bookmark}
\IfFileExists{xurl.sty}{\usepackage{xurl}}{} % add URL line breaks if available
\urlstyle{same}
\hypersetup{
  colorlinks=true,
  linkcolor={Maroon},
  filecolor={Maroon},
  citecolor={Blue},
  urlcolor={Blue},
  pdfcreator={LaTeX via pandoc}}

\author{}
\date{}

\begin{document}

{
\setcounter{tocdepth}{3}
\tableofcontents
}
\section{How it works: EEG-to-image
decoding, end to
end}\label{how-it-works-eeg-to-image-decoding-end-to-end}

A guided tour of what this repo does and
why, written for someone who wants to
understand it deeply. Read it top to
bottom once; the diagrams are a map you
can come back to. For the layer-by-layer
details, the full experiment log and the
paper map, see
\href{architecture_deep_dive.pdf}{architecture\_deep\_dive.md}.

\begin{center}\rule{0.5\linewidth}{0.5pt}\end{center}

\subsection{1. The problem in one
sentence}\label{the-problem-in-one-sentence}

A person looks at a photo while we
record their EEG for one second. From
that EEG alone, we must say
\textbf{which photo they saw}, out of
200 photos the model has never seen
during training.

We don\textquotesingle t classify into
fixed labels: the test images are of
\emph{new} concepts. Instead, we
translate brain activity into the
\textbf{same vector space as the images}
and then look up the nearest image. That
is called \textbf{zero-shot retrieval}.

\begin{figure}[H]\centering
\begin{adjustbox}{max width=\linewidth, max totalheight=0.8\textheight}
\begin{tikzpicture}[>={Stealth[length=2mm]}, box/.style={draw, rounded corners=2pt, align=center, font=\footnotesize, inner sep=4pt, fill=black!4}, every edge quotes/.style={font=\scriptsize, fill=white, inner sep=1pt, align=center, auto}]
\graph [layered layout, grow=down, level sep=16mm, sibling sep=32mm, nodes={box}, edges={semithick}] {
  EEG [as={\textbf{Brain side} \\ EEG window \\ 63 channels × 120 samples \\ (−0.2 s → +0.8 s around image onset)}];
  ENC [as={Our EEG encoder \\ (trained by us)}];
  P [as={Predicted embedding \\ 1536 numbers, length 1}];
  IMG [as={\textbf{Image side} \\ The photo shown}];
  DINO [as={DINOv2-giant \\ (frozen, by Meta AI)}];
  T [as={Image embedding \\ 1536 numbers}];
  G [as={Gallery of 200 \\ unseen images}];
  EEG ->[] ENC;
  ENC ->[] P;
  IMG ->[] DINO;
  DINO ->[] T;
  P --[dashed, "{training: pull together \\ (contrastive loss)}"] T;
  P ->[very thick, "{test: nearest of 200 \\ (cosine similarity)}"] G;
};
\end{tikzpicture}
\end{adjustbox}
\end{figure}

\begin{center}\rule{0.5\linewidth}{0.5pt}\end{center}

\subsection{2. The data:
THINGS-EEG2}\label{the-data-things-eeg2}

\begin{itemize}
\tightlist
\item
  \textbf{THINGS} is a database of 1,854
  everyday object \emph{concepts}
  (aardvark, abacus, \ldots, zucchini),
  with many natural photos per concept.
\item
  \textbf{THINGS-EEG2} (Gifford et al.,
  2022) recorded \textbf{10
  participants} watching those photos in
  a \textbf{rapid serial visual
  presentation} (RSVP): one image every
  200 ms, with a 63-channel EEG cap.

  \begin{itemize}
  \tightlist
  \item
    \textbf{Training set:} 1,654
    concepts × 10 images = 16,540
    images, each shown \textbf{4 times}
    per participant.
  \item
    \textbf{Test set:} 200 \emph{other}
    concepts × 1 image, each shown
    \textbf{80 times} per participant.
  \end{itemize}
\item
  The competition pipeline (NeuralBench)
  cuts the continuous recording into one
  \textbf{window} per image
  presentation: from 0.2 s before onset
  to 0.8 s after, filtered,
  baseline-corrected, robust-scaled,
  clipped, and \textbf{resampled to 120
  Hz}. So each example is a
  \texttt{63\ ×\ 120} matrix.
\item
  Window counts in our runs: 668,400
  training windows (train +
  NeuralBench\textquotesingle s val
  split, all concepts) and 162,240 test
  windows.
\end{itemize}

\textbf{Why so many repetitions?} A
single EEG trial is extremely noisy: the
brain response to an image is small
compared to everything else going on
(other neural activity, eye movements,
muscles, line noise). Averaging repeats
cancels the noise and keeps the
image-locked response. This one idea
comes back three times in this project:
the metric (§5), repetition-averaged
training (\#17), and ensembling.

\subsubsection{Visual processing in the
EEG, in
brief}\label{visual-processing-in-the-eeg-in-brief}

After an image appears, the visual
cortex responds within
\textasciitilde100 ms (low-level
features: edges, contrast, colour), and
by \textasciitilde150--300 ms higher
areas encode object category and
meaning. Our window (−0.2 → +0.8 s)
captures all of it. The occipital and
parietal electrodes (back of the head)
carry most of the information, which is
why the encoder learns its own
\emph{spatial} filters across channels.

\begin{center}\rule{0.5\linewidth}{0.5pt}\end{center}

\subsection{3. The target: DINOv2 image
embeddings}\label{the-target-dinov2-image-embeddings}

\href{https://arxiv.org/abs/2304.07193}{DINOv2}
is a vision model trained by
self-supervision on 142M images. Its
"giant" version turns any photo into a
\textbf{1536-d vector} in which similar
images (same object, similar shape or
texture) are close together.

We don\textquotesingle t train it.
It\textquotesingle s a fixed "coordinate
system for images". Our job is only to
map EEG into that coordinate system.
Because DINOv2 places \emph{any} image,
including the 200 unseen test images, in
that space, a good EEG→DINOv2 mapping
transfers to new concepts.
That\textquotesingle s what makes
zero-shot possible.

The embeddings are computed once and
cached (\texttt{cache/});
\texttt{scripts/mps\_patch.py} computes
them on the Mac\textquotesingle s GPU.

\begin{center}\rule{0.5\linewidth}{0.5pt}\end{center}

\subsection{4. The model and how it
learns}\label{the-model-and-how-it-learns}

\subsubsection{\texorpdfstring{4.1 The
encoder
(\texttt{submission/submission.py},
\texttt{EEGToImageEncoder})}{4.1 The encoder (submission/submission.py, EEGToImageEncoder)}}\label{the-encoder-submissionsubmission.py-eegtoimageencoder}

A small convolutional network in the
\textbf{EEGNet} family (Lawhern et al.,
2018), about 1.4M parameters:

\begin{figure}[H]\centering
\begin{adjustbox}{max width=\linewidth, max totalheight=0.8\textheight}
\begin{tikzpicture}[>={Stealth[length=2mm]}, box/.style={draw, rounded corners=2pt, align=center, font=\footnotesize, inner sep=4pt, fill=black!4}, every edge quotes/.style={font=\scriptsize, fill=white, inner sep=1pt, align=center, auto}]
\graph [layered layout, grow=down, level sep=9mm, sibling sep=6mm, nodes={box}, edges={semithick}] {
  X [as={EEG window \\ 1 × 63 channels × 120 time samples}];
  TC [as={Temporal convolution \\ 32 filters, 31 taps (\textasciitilde{}0.26 s) \\ = learned band-pass / waveform detectors}];
  SC [as={Spatial convolution \\ 64 filters spanning all 63 channels \\ = learned electrode weightings (like ICA / beamformers)}];
  P1 [as={ELU · average-pool ×4 in time · dropout}];
  SEP [as={Depthwise-separable convolution \\ 64 → 128 feature maps, 15 taps}];
  P2 [as={ELU · average-pool ×4 · dropout}];
  F [as={Flatten (128 × 7 = 896 features)}];
  H [as={Projection head \\ Linear 896→512 · BatchNorm · GELU · dropout · Linear 512→1536}];
  N [as={L2-normalise → unit-length 1536-d prediction}];
  X ->[] TC;
  TC ->[] SC;
  SC ->[] P1;
  P1 ->[] SEP;
  SEP ->[] P2;
  P2 ->[] F;
  F ->[] H;
  H ->[] N;
};
\end{tikzpicture}
\end{adjustbox}
\end{figure}

\begin{itemize}
\tightlist
\item
  \textbf{Temporal then spatial
  convolutions} mirror how EEG is
  analysed by hand: first filter in time
  (which rhythms and waveforms), then
  combine electrodes (where on the
  scalp).
\item
  \textbf{Pooling} reduces 120 time
  samples to 7: the response is slow, so
  fine timing matters less than the
  overall shape.
\item
  \textbf{The head} maps EEG features
  into DINOv2\textquotesingle s 1536-d
  space.
\item
  \textbf{Normalisation} puts every
  prediction on the unit sphere, so only
  its direction matters, matching the
  cosine-similarity metric.
\end{itemize}

Issue \#27 makes the widths configurable
(\texttt{temporal\_filters},
\texttt{spatial\_filters}, \ldots, plus
ATM-style residual blocks) to test
whether the model is too small.

\subsubsection{4.2 The loss: contrastive
learning (CLIP /
InfoNCE)}\label{the-loss-contrastive-learning-clip-infonce}

The naive approach would be to minimise
the distance between prediction and
target (regression). It works poorly
here: the EEG can\textquotesingle t
predict all 1536 numbers, and the model
learns to output a blurry "average
image". What we actually need is that
\textbf{the right image ranks above the
wrong ones}. Contrastive learning
optimises exactly that:

\begin{figure}[H]\centering
\begin{adjustbox}{max width=\linewidth, max totalheight=0.8\textheight}
\begin{tikzpicture}[>={Stealth[length=2mm]}, box/.style={draw, rounded corners=2pt, align=center, font=\footnotesize, inner sep=4pt, fill=black!4}, every edge quotes/.style={font=\scriptsize, fill=white, inner sep=1pt, align=center, auto}]
\graph [layered layout, grow=down, level sep=9mm, sibling sep=6mm, nodes={box}, edges={semithick}] {
  E [as={\textbf{One batch of 256 (EEG, image) pairs} \\ 256 predictions \\ (from EEG)}];
  I [as={256 targets \\ (DINOv2 of the shown images)}];
  S [as={256 × 256 similarity matrix \\ s\_ij = cos(pred\_i, target\_j)}];
  D [as={Divide by temperature τ}];
  CE [as={Cross-entropy per row: \\ 'which of the 256 images did EEG i see?' \\ correct answer = the diagonal}];
  L [as={Loss ↓ = right image pulled closer, \\ 255 wrong ones pushed away}];
  E ->[] S;
  I ->[] S;
  S ->[] D;
  D ->[] CE;
  CE ->[] L;
};
\end{tikzpicture}
\end{adjustbox}
\end{figure}

This is the \textbf{InfoNCE} loss (van
den Oord et al., 2018), popularised by
\textbf{CLIP} (Radford et al., 2021) for
images and text, and used for brain
decoding by Défossez et al. (2023) and
Benchetrit et al. (2024) at Meta, and by
NICE / ATM on THINGS-EEG2.

\textbf{The temperature τ} controls how
sharp the softmax is. Small τ means the
loss focuses on the hardest wrong images
(the ones that look most like the right
one). CLIP \emph{learns} τ, with a floor
at 0.01. Learning it instead of fixing
it at 0.07 was our biggest single gain
(\textbf{+10 points} held-out). In R6, τ
fell to the 0.01 floor by epoch 15, so
R15 tried a lower floor
(\texttt{min\_temperature=0.005}): τ
went lower, but the held-out score
didn\textquotesingle t improve, so the
floor stays at 0.01.

\textbf{Batch size matters} because the
other 255 items are the "wrong answers":
more negatives make a harder, more
informative task. That\textquotesingle s
one of the \#18 experiments (batch 512 /
1024).

\subsubsection{\texorpdfstring{4.3 The
optimiser and training loop
(\texttt{Solver.fit})}{4.3 The optimiser and training loop (Solver.fit)}}\label{the-optimiser-and-training-loop-solver.fit}

\begin{itemize}
\tightlist
\item
  \textbf{AdamW} (Loshchilov \& Hutter,
  2019) with weight decay 1e-4, learning
  rate 1e-3, gradient clipping at 1.0.
\item
  \textbf{Learning-rate schedule}
  (\#18): constant by default;
  \texttt{lr\_schedule=\textquotesingle{}cosine\textquotesingle{}}
  lowers it smoothly to 0 (Loshchilov \&
  Hutter, 2017), with an optional linear
  \textbf{warm-up}. Warm-up matters for
  large batches (Goyal et al., 2017).
\item
  \textbf{Per-epoch held-out scoring +
  early stopping} (\#15): after each
  epoch we score the held-out split with
  the official metric, keep the best
  epoch\textquotesingle s weights, and
  stop after \texttt{patience} epochs
  without improvement.
\item
  \textbf{Repetition averaging} (\#17,
  \texttt{-\/-avg-reps\ K}): train on
  the mean of K repeats of the same
  (participant, image) instead of single
  noisy trials.
\end{itemize}

\begin{center}\rule{0.5\linewidth}{0.5pt}\end{center}

\subsection{5. How we\textquotesingle re
scored}\label{how-were-scored}

\begin{figure}[H]\centering
\begin{adjustbox}{max width=\linewidth, max totalheight=0.8\textheight}
\begin{tikzpicture}[>={Stealth[length=2mm]}, box/.style={draw, rounded corners=2pt, align=center, font=\footnotesize, inner sep=4pt, fill=black!4}, every edge quotes/.style={font=\scriptsize, fill=white, inner sep=1pt, align=center, auto}]
\graph [layered layout, grow=down, level sep=9mm, sibling sep=6mm, nodes={box}, edges={semithick}] {
  W [as={162,240 test windows \\ 10 participants × 200 images × \textasciitilde{}80 repeats}];
  PR [as={Encoder: one 1536-d prediction per window}];
  POOL [as={Pooled metric: \\ each window ranked alone \\ against the 200 image embeddings}];
  AGG [as={Subject-aggregated metric: \\ average the \textasciitilde{}80 predictions of each \\ (participant, image) first → 2,000 queries}];
  T5a [as={top-5 accuracy: is the right image \\ among the 5 most similar? \\ (ours: 0.31)}];
  T5b [as={top-5 subject-agg \\ (ours: 0.74; top-1: 0.40) \\ = the leaderboard metric}];
  W ->[] PR;
  PR ->[] POOL;
  PR ->[] AGG;
  POOL ->[] T5a;
  AGG ->[] T5b;
};
\end{tikzpicture}
\end{adjustbox}
\end{figure}

\begin{itemize}
\tightlist
\item
  \textbf{Top-k accuracy:} for each
  query, rank the 200 gallery images by
  cosine similarity; a hit if the right
  one is in the top k. Chance level is
  5/200 = \textbf{2.5\%} for top-5 and
  0.5\% for top-1.
\item
  \textbf{Subject-aggregated} averages a
  participant\textquotesingle s repeated
  responses to the same image before
  ranking: the same noise-cancelling
  trick, applied at test time.
  It\textquotesingle s why 0.74 is so
  much higher than the pooled 0.31.
\end{itemize}

\begin{center}\rule{0.5\linewidth}{0.5pt}\end{center}

\subsection{6. How we choose models
without
cheating}\label{how-we-choose-models-without-cheating}

The warm-up leaderboard is scored on the
\textbf{public THINGS-EEG2 test set}.
Choosing models by their test score
would quietly overfit to it (teams at
1.0 on the leaderboard are almost
certainly leaking test data), and the
final phase uses different data anyway.

So we built our own \textbf{held-out
split that mimics the test set}
(\texttt{src/data/heldout.py},
\texttt{configs/heldout\_concepts.json}):

\begin{itemize}
\tightlist
\item
  200 \emph{training} concepts are
  removed entirely from training (no
  leakage of the concept);
\item
  one image per held-out concept forms a
  200-image gallery, as in the test set.
\end{itemize}

Every decision (temperature, data,
epochs, architecture) is made on this
held-out score. The test set is only
looked at for the final model, and it
has matched Codabench to two decimals on
every upload so far.

\begin{figure}[H]\centering
\begin{adjustbox}{max width=\linewidth, max totalheight=0.8\textheight}
\begin{tikzpicture}[>={Stealth[length=2mm]}, box/.style={draw, rounded corners=2pt, align=center, font=\footnotesize, inner sep=4pt, fill=black!4}, every edge quotes/.style={font=\scriptsize, fill=white, inner sep=1pt, align=center, auto}]
\graph [layered layout, grow=down, level sep=9mm, sibling sep=6mm, nodes={box}, edges={semithick}] {
  R [as={\textbf{1 · Explore (held-out)} \\ Run variants \\ (R6, R15, R7…R14) \\ seed 1, 1,454 concepts}];
  HS [as={Held-out top-5 \\ per epoch}];
  PICK [as={Keep changes worth \\ ≥ 1.5 points}];
  RT [as={\textbf{2 · Final models} \\ Retrain the winning recipe \\ on ALL concepts + val data \\ 3 seeds}];
  ENS [as={Ensemble: average the \\ 3 normalised predictions}];
  TEST [as={Score once on local test \\ (= Codabench, to 2 decimals)}];
  ZIP [as={\textbf{3 · Ship} \\ make zip-submission \\ submission.py + weights\_*.pt}];
  CB [as={Codabench runs \\ load\_model → predict}];
  RES [as={Log in RESULTS.md}];
  R ->[] HS;
  HS ->[] PICK;
  RT ->[] ENS;
  ENS ->[] TEST;
  ZIP ->[] CB;
  CB ->[] RES;
  PICK ->[] RT;
  TEST ->[] ZIP;
};
\end{tikzpicture}
\end{adjustbox}
\end{figure}

\textbf{Why retrain on all concepts?}
More concepts mean better coverage of
DINOv2\textquotesingle s space. The gaps
have been consistent: the same model
scores 13--16 points higher on test than
on held-out (held-out images were shown
\textasciitilde4 times per participant,
test images \textasciitilde80, so
subject aggregation removes more noise
on test), and the all-concept 3-seed
ensemble added about +3 more (R5: 0.5525
held-out → 0.711 test → 0.738 final).
That\textquotesingle s how we forecast:
0.59 held-out (R6) should land around
\textbf{0.75--0.78} as a final ensemble.

\textbf{Why ensembles help:} different
seeds make different errors; averaging
their (unit-length) predictions cancels
part of the error.
That\textquotesingle s the same
repetition logic again, applied to
models.

\begin{center}\rule{0.5\linewidth}{0.5pt}\end{center}

\subsection{7. The submission
contract}\label{the-submission-contract}

Codabench is \textbf{inference-only}: it
never trains. We upload a ZIP containing
\texttt{submission.py} and the weights.
The platform calls
\texttt{Solver.load\_model(meta)}
(rebuild the network for
\texttt{meta{[}"n\_chans"{]}} channels,
load the weights from a read-only
folder), then \texttt{model.predict(X)}
on each test batch. Since \#27,
\texttt{load\_model} reads the layer
sizes from the weights themselves, so
the same \texttt{submission.py} works
for any encoder width or ensemble.

\texttt{fit} and \texttt{save\_model}
are only used locally, by
\texttt{scripts/train\_solver.py}.

\begin{center}\rule{0.5\linewidth}{0.5pt}\end{center}

\subsection{8. Where we are (Oct
22)}\label{where-we-are-oct-22}

\subsubsection{Warm-up leaderboard
(THINGS-EEG2: same 10 people, 63
channels)}\label{warm-up-leaderboard-things-eeg2-same-10-people-63-channels}

{\def\LTcaptype{none} % do not increment counter
\begin{longtable}[]{@{}
  >{\raggedright\arraybackslash}p{(\linewidth - 6\tabcolsep) * \real{0.0698}}
  >{\raggedright\arraybackslash}p{(\linewidth - 6\tabcolsep) * \real{0.4884}}
  >{\raggedright\arraybackslash}p{(\linewidth - 6\tabcolsep) * \real{0.0814}}
  >{\raggedright\arraybackslash}p{(\linewidth - 6\tabcolsep) * \real{0.3605}}@{}}
\toprule\noalign{}
\begin{minipage}[b]{\linewidth}\raggedright
Date
\end{minipage} &
\begin{minipage}[b]{\linewidth}\raggedright
Change
\end{minipage} &
\begin{minipage}[b]{\linewidth}\raggedright
Held-out top-5
\end{minipage} &
\begin{minipage}[b]{\linewidth}\raggedright
Codabench top-5 (subject-agg)
\end{minipage} \\
\midrule\noalign{}
\endhead
\bottomrule\noalign{}
\endlastfoot
Oct 18 & NeuralBench EEGNet baseline &
-- & 0.29 \\
Oct 19 & Our encoder + CLIP loss,
ensemble & -- & 0.54--0.55 \\
Oct 19 & \textbf{Learnable temperature}
& 0.5175 & 0.65 \\
Oct 19 & \textbf{+25\% data}
(NeuralBench\textquotesingle s val
split) & 0.5525 & 0.71 \\
Oct 20 & Final: all concepts, 3-seed
ensemble & -- & 0.74 (top-1 0.40) \\
Oct 20 & R6: \textbf{40 epochs} instead
of 20 & 0.590 & \\
Oct 21 & \textbf{Final: R6 recipe,
3-seed ensemble} & -- & \textbf{0.78
(top-1 0.44)} \\
Oct 21 & R12: \textbf{wider
convolutions} (64/128/256 filters) &
\textbf{0.615} & \\
Oct 22 & \textbf{Final: R12 recipe
(wider convolutions), 3-seed ensemble} &
-- & \textbf{local 0.8005 (top-1
0.4655)}; upload pending a Codabench
outage \\
\end{longtable}
}

No gain, on held-out concepts:
temperature floor 0.005 (R15), cosine
schedule (R9, 0.570), batch 512 / 1024
with cosine (R10 0.581, R11 0.586:
nearly R6\textquotesingle s score in 62
instead of 105 minutes), repetition
averaging (R7/R8, cut short by early
stopping, to re-test), a wider head only
(R13, trailing, cancelled).

\subsubsection{The final phase: unseen
people, 32
channels}\label{the-final-phase-unseen-people-32-channels}

The sealed final is a private
\textbf{Alljoined} cohort: \textbf{11
new participants, 32-channel Emotiv
headset}, new images. Updated schedule
(organisers, Oct 22): warm-up and
registration close \textbf{Nov 1}; the
sealed phase runs \textbf{Nov 5 -- Nov
21} with only \textbf{5 sealed uploads
per participant in total}, so local
validation decides what we upload. We
now measure exactly the
final\textquotesingle s setting:

\begin{itemize}
\tightlist
\item
  \textbf{THINGS-EEG2 held-out
  participants}
  (\texttt{-\/-holdout-subjects},
  participants 6 and 7 never seen, 10
  galleries): the R6 recipe scores
  \textbf{0.233} on unseen people,
  against \textbf{0.635} when the same
  two people were in training. Unseen
  people are the main problem.
\item
  \textbf{Alljoined-1.6M} (downloaded
  Oct 21: 20 participants, 32 channels,
  \textasciitilde1.66M windows, exactly
  THINGS-EEG2\textquotesingle s images):
  trained on 16 participants, scored on
  4 unseen ones over the 200 test
  images: \textbf{≈ 0.27} (A1).
\end{itemize}

What helps on unseen people
(THINGS-EEG2, one change each from
0.233):

{\def\LTcaptype{none} % do not increment counter
\begin{longtable}[]{@{}
  >{\raggedright\arraybackslash}p{(\linewidth - 2\tabcolsep) * \real{0.6393}}
  >{\raggedright\arraybackslash}p{(\linewidth - 2\tabcolsep) * \real{0.3607}}@{}}
\toprule\noalign{}
\begin{minipage}[b]{\linewidth}\raggedright
Change
\end{minipage} &
\begin{minipage}[b]{\linewidth}\raggedright
Unseen-people score
\end{minipage} \\
\midrule\noalign{}
\endhead
\bottomrule\noalign{}
\endlastfoot
\textbf{Per-window normalisation}
(\texttt{window\_norm}) & \textbf{0.274}
(+4.1) \\
Wider convolutions & 0.253 (+2.0, 3.5×
slower) \\
Spatial attention (montage-independent
input layer) & 0.255 (+1.6 on
average) \\
Augmentation (time shift, channel
dropout) & 0.230 (none) \\
Participant-adversarial training & 0.190
(hurts) \\
\textbf{Fewer training participants:} 2
/ 4 / 6 / 8 & 0.161 / 0.183 / 0.187 /
0.233 \\
\end{longtable}
}

\textbf{More participants is the biggest
lever within one dataset.} Two ways to
get more people, tried on
Alljoined\textquotesingle s unseen
participants (A1, from scratch: ≈
0.265--0.278):

{\def\LTcaptype{none} % do not increment counter
\begin{longtable}[]{@{}
  >{\raggedright\arraybackslash}p{(\linewidth - 4\tabcolsep) * \real{0.0882}}
  >{\raggedright\arraybackslash}p{(\linewidth - 4\tabcolsep) * \real{0.7353}}
  >{\raggedright\arraybackslash}p{(\linewidth - 4\tabcolsep) * \real{0.1765}}@{}}
\toprule\noalign{}
\begin{minipage}[b]{\linewidth}\raggedright
Run
\end{minipage} &
\begin{minipage}[b]{\linewidth}\raggedright
Change
\end{minipage} &
\begin{minipage}[b]{\linewidth}\raggedright
Unseen-people score
\end{minipage} \\
\midrule\noalign{}
\endhead
\bottomrule\noalign{}
\endlastfoot
A4 & joint training: +
THINGS-EEG2\textquotesingle s 10
participants (reduced to
Alljoined\textquotesingle s 32 channels)
& 0.241 \\
A3 & frozen \textbf{REVE} (an EEG model
pretrained on many people) +
attention-pooling head & 0.245 \\
\end{longtable}
}

Neither beats Alljoined alone. Oct 22
re-tests with smoothed early stopping
(best 3-epoch mean):

{\def\LTcaptype{none} % do not increment counter
\begin{longtable}[]{@{}
  >{\raggedright\arraybackslash}p{(\linewidth - 4\tabcolsep) * \real{0.0811}}
  >{\raggedright\arraybackslash}p{(\linewidth - 4\tabcolsep) * \real{0.5000}}
  >{\raggedright\arraybackslash}p{(\linewidth - 4\tabcolsep) * \real{0.4189}}@{}}
\toprule\noalign{}
\begin{minipage}[b]{\linewidth}\raggedright
Run
\end{minipage} &
\begin{minipage}[b]{\linewidth}\raggedright
Change
\end{minipage} &
\begin{minipage}[b]{\linewidth}\raggedright
Unseen-people score (smoothed)
\end{minipage} \\
\midrule\noalign{}
\endhead
\bottomrule\noalign{}
\endlastfoot
B1 & Alljoined only (16 participants) &
\textbf{0.2725} \\
B2 & + per-window normalisation & 0.2696
(no gain here, unlike
THINGS-EEG2\textquotesingle s +4.1) \\
B3 & + THINGS-EEG2 joint training +
per-window normalisation & 0.2471 \\
\end{longtable}
}

\textbf{THINGS-EEG2
doesn\textquotesingle t help Alljoined}:
an Alljoined model even scores only
0.066 on THINGS-EEG2\textquotesingle s
test (chance 0.025). The two datasets
differ too much (lab-grade cap vs
Emotiv, 200 vs 100 ms presentations). So
the final model will be trained on
Alljoined.

\textbf{REVE, fine-tuned, as an ensemble
member.} NeuralBench\textquotesingle s
reference (84.75\% top-5 on THINGS-EEG2,
above our 0.80) \textbf{fine-tunes
REVE}; frozen, it gave 0.245 on
Alljoined\textquotesingle s unseen
people. Fine-tuning its top 4 layers
gives 0.259 alone, but \textbf{averaged
with two of our encoders: 0.3075},
against 0.2737 for our best single
encoder, our first score above 0.30 on
unseen people.

\textbf{Two-stage training is the best
single model:} pretrain on THINGS-EEG2 +
Alljoined, then fine-tune on Alljoined
alone: \textbf{0.308 / 0.311 / 0.316}
for three seeds (smoothed), +3.6 to +4.3
over Alljoined alone. \textbf{Best
ensemble: three two-stage seeds + one
fine-tuned REVE ≈ 0.335} (all 63
combinations of six models scored; REVE
should get about a quarter of the
weight). Target for Nov 5: ≥ 0.35.

\textbf{Cross-validation (Oct 25--26)
corrects the level:} our 4 validation
people turned out easier than average.
Over 5 folds covering all 20
participants, the final recipe
(two-stage + REVE) scores \textbf{0.240}
on average unseen people (+2.5 over
single-stage, better in 4 of 5 folds),
single-stage 0.215; folds range from
0.17 to 0.32 depending on the people.
Realistic sealed expectation:
\textbf{\textasciitilde0.24--0.26
top-5}. Details:
\href{architecture_deep_dive.pdf}{architecture\_deep\_dive.md}
§13.

\textbf{Final models (Oct 23):} the
recipe retrained on all 20 Alljoined
participants with fixed epoch counts
(three two-stage seeds + REVE), packaged
as one Codabench submission (87 MB,
contract check passed): our first
sealed-phase candidate. Full details:
\href{architecture_deep_dive.pdf}{architecture\_deep\_dive.md}
§12. The whole REVE story, from its
architecture to our fine-tuning
implementation, with diagrams and a
slide outline:
\href{architecture_deep_dive.pdf}{architecture\_deep\_dive.md}
§11. Details in
\href{architecture_deep_dive.pdf}{architecture\_deep\_dive.md}
§9--§10.

\begin{center}\rule{0.5\linewidth}{0.5pt}\end{center}

\subsection{9. Glossary}\label{glossary}

{\def\LTcaptype{none} % do not increment counter
\begin{longtable}[]{@{}
  >{\raggedright\arraybackslash}p{(\linewidth - 2\tabcolsep) * \real{0.2958}}
  >{\raggedright\arraybackslash}p{(\linewidth - 2\tabcolsep) * \real{0.7042}}@{}}
\toprule\noalign{}
\begin{minipage}[b]{\linewidth}\raggedright
Term
\end{minipage} &
\begin{minipage}[b]{\linewidth}\raggedright
Meaning
\end{minipage} \\
\midrule\noalign{}
\endhead
\bottomrule\noalign{}
\endlastfoot
\textbf{Epoch / window / trial} & One
image presentation\textquotesingle s
slice of EEG (\texttt{63\ ×\ 120} here).
\emph{Training} epoch = one pass over
all windows; context tells which \\
\textbf{Embedding} & A vector
representing an item; similar items have
nearby vectors \\
\textbf{Cosine similarity} & Dot product
of unit vectors: 1 = same direction, 0 =
unrelated \\
\textbf{Zero-shot} & Recognising classes
never seen in training, via a shared
embedding space \\
\textbf{Contrastive loss / InfoNCE} &
"Pick the right partner among the batch"
cross-entropy \\
\textbf{Temperature τ} & Divides
similarities before the softmax; smaller
= sharper \\
\textbf{Held-out split} & Our private
test-like validation set (200 training
concepts) \\
\textbf{Subject-aggregated} & Average a
participant\textquotesingle s repeats of
an image before scoring \\
\textbf{Seed} & The random-number start;
different seeds give different models
from the same recipe \\
\textbf{Early stopping} & Keep the best
epoch\textquotesingle s weights, stop
when the validation score stops
improving \\
\textbf{RSVP} & Rapid serial visual
presentation: images flashed in quick
succession \\
\end{longtable}
}

\begin{center}\rule{0.5\linewidth}{0.5pt}\end{center}

\subsection{10. Reading
list}\label{reading-list}

The full, checked bibliography (with
which references were verified) is in
\href{architecture_deep_dive.pdf\#7-paper-map-and-bibliography}{architecture\_deep\_dive.md
§7}. The arXiv numbers for CLIP, NICE,
ATM, Benchetrit et al. and the UBP
details could not be verified offline;
search the title if a link
doesn\textquotesingle t resolve.

\subsubsection{Start here: the core
ideas (≈ 1
evening)}\label{start-here-the-core-ideas-1-evening}

\begin{enumerate}
\def\labelenumi{\arabic{enumi}.}
\tightlist
\item
  \textbf{Lilian Weng, "Contrastive
  Representation Learning"} (blog,
  2021). The clearest introduction to
  InfoNCE, negatives and temperature.
\item
  \textbf{Radford et al., "Learning
  Transferable Visual Models From
  Natural Language Supervision" (CLIP,
  2021)}, arXiv:2103.00020. Read §2.3
  and Figure 3 (the pseudo-code is our
  loss, learnable temperature included).
\item
  \textbf{Song et al., "Decoding Natural
  Images from EEG for Object
  Recognition" (NICE, ICLR 2024)},
  arXiv:2308.13234. \emph{The} paper for
  our exact setup: THINGS-EEG2, 200-way
  zero-shot, contrastive EEG↔image
  training. Its results tables are the
  yardstick.
\end{enumerate}

\subsubsection{Our task
specifically}\label{our-task-specifically}

\begin{enumerate}
\def\labelenumi{\arabic{enumi}.}
\setcounter{enumi}{3}
\tightlist
\item
  \textbf{Gifford et al., "A large and
  rich EEG dataset for modeling human
  visual object recognition"
  (NeuroImage, 2022)},
  \url{doi:10.1016/j.neuroimage.2022.119754}.
  The THINGS-EEG2 dataset paper:
  paradigm, preprocessing.
\item
  \textbf{Hebart et al., "THINGS: A
  database of 1,854 object concepts and
  more than 26,000 naturalistic object
  images" (PLoS ONE, 2019).}
\item
  \textbf{Li et al., "Visual Decoding
  and Reconstruction via EEG Embeddings
  with Guided Diffusion" (ATM, NeurIPS
  2024)}, arXiv:2403.07721. A stronger
  encoder (issue \#12) and the residual
  projection head behind
  \texttt{head\_blocks}.
\item
  \textbf{"Bridging the Vision-Brain Gap
  with an Uncertainty-Aware Blur Prior"
  (UBP, CVPR 2025).} Among the best
  published THINGS-EEG2 results, and the
  source of issue \#26\textquotesingle s
  blurred-image targets.
\item
  \textbf{The Alljoined papers (Xu et
  al.)}: "Alljoined --- A dataset for
  EEG-to-Image decoding" (2024,
  arXiv:2404.05553) and "Alljoined-1.6M:
  A Million-Trial EEG-Image Dataset for
  Evaluating Affordable Brain-Computer
  Interfaces" (2025, arXiv:2508.18571,
  Emotiv headset). The final
  phase\textquotesingle s data family;
  read before M4.
\end{enumerate}

\subsubsection{Brain decoding at Meta
(same recipe, other
modalities)}\label{brain-decoding-at-meta-same-recipe-other-modalities}

\begin{enumerate}
\def\labelenumi{\arabic{enumi}.}
\setcounter{enumi}{8}
\tightlist
\item
  \textbf{Défossez et al., "Decoding
  speech perception from non-invasive
  brain recordings" (Nature Machine
  Intelligence, 2023)},
  arXiv:2208.12266. The CLIP loss for
  MEG/EEG, and the subject layer that
  handles different people.
\item
  \textbf{Benchetrit, Banville \& King,
  "Brain decoding: toward real-time
  reconstruction of visual perception"
  (ICLR 2024)}, arXiv:2310.19812.
  MEG→image retrieval with DINOv2-style
  targets; very close to what we do.
\end{enumerate}

\subsubsection{The building
blocks}\label{the-building-blocks}

\begin{enumerate}
\def\labelenumi{\arabic{enumi}.}
\setcounter{enumi}{10}
\tightlist
\item
  \textbf{Lawhern et al., "EEGNet: a
  compact convolutional neural network
  for EEG-based brain--computer
  interfaces" (J. Neural Engineering,
  2018)}, arXiv:1611.08024. Our
  encoder\textquotesingle s ancestor;
  explains temporal → spatial →
  separable convolutions.
\item
  \textbf{Oquab et al., "DINOv2:
  Learning Robust Visual Features
  without Supervision" (2023)},
  arXiv:2304.07193. What our target
  space is.
\item
  \textbf{van den Oord et al.,
  "Representation Learning with
  Contrastive Predictive Coding"
  (2018)}, arXiv:1807.03748. Where
  InfoNCE comes from.
\item
  \textbf{Loshchilov \& Hutter,
  "Decoupled Weight Decay
  Regularization" (AdamW, ICLR 2019)},
  arXiv:1711.05101, and \textbf{"SGDR:
  Stochastic Gradient Descent with Warm
  Restarts" (ICLR 2017)},
  arXiv:1608.03983 (cosine schedule).
\item
  \textbf{Goyal et al., "Accurate, Large
  Minibatch SGD: Training ImageNet in 1
  Hour" (2017)}, arXiv:1706.02677.
  Linear learning-rate scaling and
  warm-up for big batches (\#18).
\end{enumerate}

\subsubsection{Generalising to new
people (added Oct 21; from memory, check
titles)}\label{generalising-to-new-people-added-oct-21-from-memory-check-titles}

\begin{itemize}
\tightlist
\item
  \textbf{Kolesnikov et al., "Big
  Transfer (BiT)" (ECCV 2020).} More
  (and more varied) pretraining data
  matters most when the test data
  differs from training: our "more
  participants" lever.
\item
  \textbf{El Ouahidi et al., REVE
  (2025).} The EEG foundation model we
  transfer from (§8, deep dive §10.4).
\item
  \textbf{Ganin \& Lempitsky (2015),
  DANN, and Tzeng et al. (2015), domain
  confusion.} Adversarial domain
  invariance, which we tried for
  participants (it hurt: deep dive
  §10.3).
\end{itemize}

\subsubsection{EEG fundamentals (if the
neuroscience side is
new)}\label{eeg-fundamentals-if-the-neuroscience-side-is-new}

\begin{enumerate}
\def\labelenumi{\arabic{enumi}.}
\setcounter{enumi}{15}
\tightlist
\item
  \textbf{MNE-Python tutorials}
  (mne.tools): "Overview of MEG/EEG
  analysis", epochs, evoked responses.
  Hands-on, with the same concepts as
  our pipeline (filtering, epoching,
  baseline correction).
\item
  \textbf{Steven Luck, \emph{An
  Introduction to the Event-Related
  Potential Technique}} (MIT Press). Why
  averaging repeats works, and what the
  visual responses (P1, N1, N170\ldots)
  are.
\item
  \textbf{Mike X Cohen, \emph{Analyzing
  Neural Time Series Data}} (MIT Press,
  2014). Filtering, time-frequency,
  spatial filters; very readable.
\end{enumerate}

\subsubsection{Suggested
order}\label{suggested-order}

Weng\textquotesingle s blog → CLIP §2.3
→ NICE → this doc\textquotesingle s
§4--6 again → EEGNet → Défossez → ATM /
UBP → Alljoined. Then read
\texttt{submission/submission.py} top to
bottom: after these, every line should
make sense.

\end{document}
